\documentclass[letter,11pt]{article}
\usepackage{hyperref}
\usepackage{graphicx}
\usepackage{fancyhdr}
\usepackage{multicol}
\pagestyle{fancy}
\usepackage[letterpaper, margin=1in]{geometry}
\geometry{letterpaper}
\usepackage{parskip} % Disable initial indent

\usepackage[utf8]{inputenc}
\fancyhf{}
\renewcommand{\headrulewidth}{0pt} % Remove default underline from header package
\rhead{CYBR/CMSC 691 - Lab 3}
%\rhead{}
\lhead{\begin{picture}(0,0) \put(0,-10){\includegraphics[width=1.1cm]{Images/UMBC-vertical}} \end{picture}}
\cfoot{\thepage}
\rfoot{Spring 2020}
\lfoot{Lab 3}
%\AtEndDocument{\vfill \hfill \LaTeX}

\begin{document}

\section*{Lab 3: Malware Networks}
\paragraph{}Provided in the VM is the example data (and code) to create malware networks as shown in Chapters 4 \& 5. Following along from the book, write code which extracts the data from the malware, builds connections, and displays the results.

\textbf{Points:} 20

\section*{Questions}
\paragraph{}In a separate document, provide your answers to the following questions. The expected format is a Word document, or similar, with a few images demonstrating your progress. The assignment will be submitted via Blackboard. You may use any tool(s) that you wish to answer the questions below.

\begin{itemize}
    \item Chapter 4
    \begin{enumerate}
        \item Show the resulting network for the APT1 malware as provided by the textbook authors.
        \item Test the code on some of the other malware data provided.
        \item Try to improve the pattern-matching function in the textbook author's code to get better domain strings. How does the network look now?
        \item What are some additional static features, besides host name, which might be useful to show how malware is (or could be) related?
        \item What are some limitations of this approach (creating networks from static features)?
        \item What are some ways where this approach could produce false positives (showing a non-existent relationship)?
    \end{enumerate}
    \item Chapter 5
    \begin{enumerate}
        \item Create a visual network using the textbook's code using any of the provided malware.
        \item Briefly describe the algorithms were used by the textbook authors.
    \end{enumerate}
    \item Optionally include goodware with your malware networks, but use a different color so it can be distinguished.
    \item Don't forget to include screenshots.
\end{itemize}
\end{document}